\section{Experiments}
\label{sec:experiments}

\subsection{Ablation: VLM Evolution Loop Effectiveness}
\label{sec:exp:ablation}

We first validate that the VLM-guided evolution loop improves rendering quality
compared to single-pass rendering without iterative refinement.
We construct two versions of the 50-object ARIS-Rotate rendering set:

\begin{itemize}
  \item \textbf{No-evolution (direct render):} The 3D mesh is inserted into the Blender scene
    and rendered at the default control state without any VLM review or parameter adjustment.
  \item \textbf{Full pipeline (VLM evolution):} The complete \aris{} loop is applied,
    with the AI agent reading VLM free-form feedback and selecting rendering actions
    until the reviewer issues a keep verdict.
\end{itemize}

\begin{table}[h]
\caption{Ablation: effect of VLM evolution loop on rendering quality.
Keep rate = fraction of objects receiving a VLM ``keep'' verdict.
Avg. hybrid score combines VLM quality ratings and CV-based metrics.}
\label{tab:ablation}
\centering
\begin{tabular}{lccc}
\toprule
Setting & Avg. Quality Score $\uparrow$ & Keep Rate (\%) $\uparrow$ & Avg. Rounds \\
\midrule
No-evolution (direct render)         & \tbd{XX.X} & \tbd{XX.X} & 0 \\
Full pipeline (VLM evolution, ours)  & \tbd{XX.X} & \tbd{XX.X} & \tbd{X.X} \\
\bottomrule
\end{tabular}
\end{table}

\begin{figure}[t]
  \centering
  % TODO: Insert ablation comparison figure
  \fbox{\parbox{0.95\linewidth}{\centering\vspace{2.5cm}
    \textbf{[Figure 4: Ablation Comparison]}\\
    3--4 representative objects.
    Left: No-evolution (direct render). Right: Full pipeline (VLM evolution).\\
    Annotate primary issue fixed (flat lighting, color shift, shadow, etc.)
    \vspace{2.5cm}}}
  \caption{
    Qualitative comparison between single-pass rendering (no evolution)
    and the full \aris{} evolution pipeline.
    The VLM loop corrects common issues including flat lighting, color shift,
    and missing contact shadows.
  }
  \label{fig:ablation}
\end{figure}


\subsection{LoRA Fine-Tuning Setup}
\label{sec:exp:lora_setup}

We evaluate the downstream training value of ARIS-Rotate by fine-tuning
Qwen Image Edit 2511~\cite{PLACEHOLDER_qwen_image_edit_verify} with LoRA
on the training split and evaluating on the test split.

\paragraph{Base model.}
Qwen Image Edit 2511 is an instruction-following image editing model
that takes a source image and a text instruction as input and produces an edited image.
We use the model's standard inference interface without modification.

\paragraph{Training configuration.}
\begin{itemize}
  \item LoRA rank: \tbd{r=XX}, alpha: \tbd{XX}, target modules: \tbd{list}
  \item Learning rate: \tbd{1e-4}, batch size: \tbd{X}, epochs: \tbd{XX}
  \item Input: source image (front view) + text instruction
  \item Target: edited image (target viewpoint)
\end{itemize}

\paragraph{Compared conditions.}
We compare three conditions:
\begin{enumerate}
  \item \textbf{Base model}: Qwen Image Edit 2511 without fine-tuning.
  \item \textbf{LoRA (no-evolution data)}: LoRA fine-tuned on data from
    single-pass rendering (no VLM refinement).
  \item \textbf{LoRA (full pipeline data)}: LoRA fine-tuned on quality-approved data
    from the full \aris{} evolution pipeline.
\end{enumerate}


\subsection{Quantitative Results}
\label{sec:exp:quant}

We evaluate on the object-disjoint test split (10 objects, 70 pairs)
using standard image quality metrics: PSNR, SSIM, and LPIPS~\cite{PLACEHOLDER_zhang2018lpips_verify}.

\begin{table}[h]
\caption{
  Quantitative evaluation of rotation editing quality on the ARIS-Rotate test set.
  $\uparrow$ = higher is better; $\downarrow$ = lower is better.
  Bold denotes best result per metric.
}
\label{tab:quant_main}
\centering
\begin{tabular}{lccc}
\toprule
Method & PSNR $\uparrow$ & SSIM $\uparrow$ & LPIPS $\downarrow$ \\
\midrule
Base model (no fine-tuning)        & \tbd{XX.X} & \tbd{0.XX} & \tbd{0.XX} \\
LoRA on no-evolution data          & \tbd{XX.X} & \tbd{0.XX} & \tbd{0.XX} \\
LoRA on full pipeline data (ours)  & \textbf{\tbd{XX.X}} & \textbf{\tbd{0.XX}} & \textbf{\tbd{0.XX}} \\
\bottomrule
\end{tabular}
\end{table}

\paragraph{Per-view breakdown.}
Table~\ref{tab:perviewresults} shows results broken down by target viewpoint.
Rotation to nearby views (front-right, front-left) is generally easier than
rotation to back views, which require hallucinating occluded geometry.

\begin{table}[h]
\caption{
  Per-viewpoint PSNR comparison between base model and LoRA on full pipeline data.
}
\label{tab:perviewresults}
\centering
\begin{tabular}{llcc}
\toprule
Target View & Azimuth & Base PSNR & LoRA PSNR \\
\midrule
front-right view & $45^\circ$  & \tbd{XX.X} & \tbd{XX.X} \\
right side view  & $90^\circ$  & \tbd{XX.X} & \tbd{XX.X} \\
back-right view  & $135^\circ$ & \tbd{XX.X} & \tbd{XX.X} \\
back view        & $180^\circ$ & \tbd{XX.X} & \tbd{XX.X} \\
back-left view   & $225^\circ$ & \tbd{XX.X} & \tbd{XX.X} \\
left side view   & $270^\circ$ & \tbd{XX.X} & \tbd{XX.X} \\
front-left view  & $315^\circ$ & \tbd{XX.X} & \tbd{XX.X} \\
\midrule
\textbf{Mean} & & \tbd{XX.X} & \tbd{XX.X} \\
\bottomrule
\end{tabular}
\end{table}


\subsection{Qualitative Results}
\label{sec:exp:qual}

\begin{figure}[t]
  \centering
  % TODO: Insert LoRA qualitative comparison figure
  \fbox{\parbox{0.95\linewidth}{\centering\vspace{3cm}
    \textbf{[Figure 5: LoRA vs Base Model Qualitative Comparison]}\\
    Rows: 3--4 representative objects.\\
    Columns: Input (front view) $|$ Base model output $|$ LoRA output $|$ Ground truth.\\
    Show difficult views (back view, left side view) where improvement is visible.
    \vspace{3cm}}}
  \caption{
    Qualitative comparison of rotation editing results.
    LoRA fine-tuned on ARIS-Rotate data produces more geometrically consistent
    target views compared to the base model, especially for large rotation angles.
  }
  \label{fig:qual_lora}
\end{figure}
